\medskip

\emph{Les probabilités demandées seront exprimées sous forme de fractions irréductibles}

\medskip

\textbf{Partie A}

\medskip

On lance trois fois de suite une pièce de monnaie bien équilibrée. On note $X$ la variable aléatoire qui compte le nombre de fois, sur les trois lancers, où la pièce est retombée du côté \og Face \fg.

\medskip

\begin{enumerate}
\item 

On peut dresser un arbre pondéré de probabilités :

\begin{center}
\pstree[treemode=R,nodesepA=0pt,nodesepB=3pt,levelsep=3cm,treesep=0.75cm]{\TR{}}
{\pstree{\TR{$P$~~}\taput{$\frac12$}}
	{\pstree{\TR{$P$~~}\taput{$\frac12$}}
		{\TR{$P$~~0}\taput{$\frac12$}
		\TR{$F$~~1}\tbput{$\frac12$}
		}
	\pstree{\TR{$F$~~}\tbput{$\frac12$}}
		{\TR{$P$~~1}\taput{$\frac12$}
		\TR{$F$~~2}\tbput{$\frac12$}
		}	
	}
\pstree{\TR{$F$~~}\tbput{$\frac12$}}
	{\pstree{\TR{$P$~~}\taput{$\frac12$}}
		{\TR{$P$~~1}\taput{$\frac12$}
		\TR{$F$~~2}\tbput{$\frac12$}
		}
	\pstree{\TR{$F$~~}\tbput{$\frac12$}}
		{\TR{$P$~~2}\taput{$\frac12$}
		\TR{$F$~~3}\tbput{$\frac12$}
		}	
	}	
}
\end{center}
\item Recopier et compléter le tableau suivant donnant la loi de probabilité de $X$
\begin{center}
\renewcommand\arraystretch{2}
\begin{tabularx}{\linewidth}{|*{5}{>{\centering \arraybackslash}X|}}\hline

$k$& 0 & 1 & 2& 3 \\ \hline
$P(X = k)$&$\dfrac18$	&$\dfrac38$	&$\dfrac38$	&$\dfrac18$\\ \hline
\end{tabularx}
\end{center}
\end{enumerate}

\medskip

\textbf{Partie B}

\medskip

Voici les règles d'un jeu où le but est d'obtenir trois pièces du côté \og Face \fg{} en un ou deux essais :

\begin{itemize}[label=$\bullet~$]
\item On lance trois pièces équilibrées :
	\begin{itemize}[label=$\circ~$]
	\item Si les trois pièces sont tombées du côté \og Face \fg, la partie est gagnée ;
	\item Sinon, les pièces tombées du côté \og Face\fg{} sont conservées et on relance
celles tombées du côté \og Pile \fg.
	\end{itemize}
\item La partie est gagnée si on obtient trois pièces du côté \og Face \fg, sinon elle est
perdue.
\end{itemize}

On considère les évènements suivants :
\begin{itemize}[label=$\bullet~$]
\item $G$ : \og la partie est gagnée \fg.

Et pour tout entier $k$ compris entre 0 et 3, les évènements:
\item $A_k$ : \og $k$ pièces sont tombées du côté \og Face\fg{} au premier lancer \fg.
\end{itemize}

\medskip

\begin{enumerate}
\item Démontrer que $P_{A_1}(G) = \dfrac14$.
\item Recopier et compléter l'arbre pondéré ci-dessous:

\begin{center}
\pstree[treemode=R,levelsep=3cm]{\TR{}}
{\pstree{\TR{$A_0~~$}\taput{$\frac18$}}
	{\TR{$G$}\taput{$\frac18$}
	\TR{$\overline{G}$}\tbput{$\frac78$}
	}
	\pstree{\TR{$A_1~~$}\taput{$\frac38$}}
	{\TR{$G$}\taput{$\frac14$}
	\TR{$\overline{G}$}\tbput{$\frac34$}
	}
	\pstree{\TR{$A_2~~$}\tbput{$\frac38$}}
	{\TR{$G$}\taput{$\frac12$}
	\TR{$\overline{G}$}\tbput{$\frac12$}
	}
	\pstree{\TR{$A_3~~$}\tbput{$\frac18$}}
	{\TR{$G$}\taput{$1$}	
	}
}
\end{center}

\item 
D'après la loi des probabilités totales on a :

$p = p(A_0) \times p_{A_0}(G) + p(A_1) \times p_{A_1}(G) + p(A_2) \times p_{A_2}(G) + p(A_3) \times p_{A_3}(G) = $

$\dfrac18 \times \dfrac18 + \dfrac38 \times \dfrac14 + \dfrac38 \times \dfrac12 + \dfrac38 \times 1 = \dfrac18 \times \dfrac18 + \dfrac38 \times \dfrac28 + \dfrac38 \times \dfrac48 + \dfrac38 \times \dfrac88 = \dfrac{1 + 6 + 12 + 8}{64} = \dfrac{27}{64}$.
\item 
Il faut trouver $p_G(A_1) = \dfrac{p\left(G \cap A_1\right)}{p(G)} = \dfrac{p\left(A_1 \cap G\right)}{p(G)} = \dfrac{\dfrac38 \times \dfrac14}{\dfrac{27}{64}} = \dfrac{\dfrac{3}{32}}{\dfrac{27}{64}} = \dfrac{3}{32} \times \dfrac{64}{27} = \dfrac29$.
\item 
Les parties sont à chaque fois indépendantes et la probabilité de gagner est à chaque fois égale à $\np{0,421875} = \dfrac{27}{64}$. Si sur $n$ parties jouées on appelle $X$ la variable aléatoire comptant le nombre de parties gagnées, $X$ suit une loi binomiale de paramètres $n$ et $p = \dfrac{27}{64}$.

On a donc $p(X = 0) = \displaystyle\binom{n}{0} \times \left(\dfrac{27}{64}\right)^0 \times \left(1 - \dfrac{27}{64}\right)^n = \left(\dfrac{37}{64}\right)^n$.

Il faut trouver $n \in \N$ tel que $p(X \geqslant 1) = 1 - \left(\dfrac{37}{64}\right)^n > 0,95 \iff \left(\dfrac{37}{64}\right)^n < 0,05 \iff n \ln \dfrac{37}{64} < \ln 0,05$ (par croissance de la fonction logarithme népérien), soit $n > \dfrac{\ln 0,05}{\ln \frac{37}{64}}$ car $\dfrac{1}{\ln \frac{37}{64}} < 0$.

Or la calculatrice donne $\dfrac{\ln 0,05}{\ln \frac{37}{64}} \approx 5,46$. Il faut donc faire au moins 6 parties.
\end{enumerate}

\bigskip
